% !TeX root = ../main.tex
\chapter*{Lời nói đầu}

\phantomsection\addcontentsline{toc}{chapter}{Lời nói đầu}

Thời đại công nghệ phát triển hiện nay, khoa học công nghệ nói chung và Robot nói riêng đang phát triển mạnh mẽ ở các nước phát triển và đang phát triển, góp phần mang lại hiệu quả kinh tế cao. Về việc nghiên cứu, chế tạo robot để ứng dụng trong đời sống và sản xuất là cần thiết.

Sự ra đời của Robot đã giúp đỡ con người trong những công việc cần độ chính xác lặp lại cao hoặc thay thế con người làm việc trong những môi trường gây hại đến sức khỏe con người như các phòng hóa chất, hầm mỏ,.... Theo dạng chuỗi động học có 2 loại Robot công nghiệp phổ biến hiện nay là Robot nối tiếp và tay máy Robot song song. Tay máy Robot song song có cấu tạo là một bàn động được liên kết với bệ cố định bằng một số chuỗi độc lập để tạo thành 1 vòng kín được áp dụng cho các tác vụ gắp, đóng gói hoặc định vị chính xác với tốc độ cao,.... Ngược lại, tay máy Robot nối tiếp ( Robot dạng chuỗi) được cấu tạo từ các khâu nối lần lượt với nhau bằng các khớp quay hoặc khớp trượt, tạo thành một chuỗi từ đế cho đến cuối khâu cuối. Robot dạng chuỗi được áp dụng chủ yếu trong các tác vụ gắp, lắp ráp, sắp xếp và đặc biệt là hàn với độ chính xác cao.

“Đồ án Thiết kế hệ thống cơ khí - Cơ điện tử thông minh” là một học phần bắt buộc trong chương trình Module 4: Hệ thống Cơ điện tử thông minh. Vì vậy trong báo cáo này, em tập trung vào việc tính toán, mô phỏng Robot chuỗi ROKAE SR5/0.9C trên thư viện ROS2 và MuJoCo. Báo cáo của em được thực hiện dưới sự hướng dẫn của TS. Trương Quốc Chiến.

Vì thời gian và kiến thức có hạn, sẽ không tránh khỏi những sai sót ngoài ý muốn. Do vậy, em mong nhận được sự góp ý và nhận xét của thầy để đề tài được hoàn thiện hơn.

Em xin chân thành cảm ơn sự hướng dẫn tận tình của TS. Trương Quốc Chiến đã giúp em hoàn thành được đề tài!

\begin{flushright}\itshape Hà Nội, ngày \ldots{} tháng \ldots{} năm 2026\end{flushright}
